\chapter{\texorpdfstring{The $\Fz$ Operator}{The Fζ Operator}}
\label{ch:fz}

\section{Difference Operators and Coefficient Uniqueness}

\begin{definition}[Symmetric difference operator]
\label{def:diff5}
\lean{FUST.DζOperator.Diff5}
\[
  \mathrm{Diff}_5[f](z) =
    f(\vp^2 z) + f(\vp z) - 4f(z) + f(\ps z) + f(\ps^2 z).
\]
\end{definition}

\begin{definition}[Antisymmetric difference operator]
\label{def:diff6}
\lean{FUST.DζOperator.Diff6}
\[
  \mathrm{Diff}_6[f](z) =
    f(\vp^3 z) - 3f(\vp^2 z) + f(\vp z) - f(\ps z) + 3f(\ps^2 z) - f(\ps^3 z).
\]
\end{definition}

\begin{theorem}[Coefficient uniqueness]
\label{thm:coeff-unique}
\lean{FUST.DζOperator.FUST_coefficient_uniqueness}
The kernel conditions
$\mathrm{Diff}_5[1]=\mathrm{Diff}_5[\mathrm{id}]=0$ uniquely determine
the coefficients $(a,b)=(-1,-4)$. Similarly,
$\mathrm{Diff}_6[\mathrm{id}]=\mathrm{Diff}_6[x^2]=0$ uniquely determine
$(A,B)=(3,1)$.
\end{theorem}

\begin{theorem}[Diff7 reduction]
\label{thm:diff7-reduction}
\lean{FUST.DζOperator.Diff7_kernel_equals_Diff6_kernel}
Any order-7 antisymmetric operator satisfying the same kernel conditions
has $\ker(\mathrm{Diff}_7) = \ker(\mathrm{Diff}_6) = \{1, z, z^2\}$.
\end{theorem}

\section{\texorpdfstring{$\Dz$ and $\Fz$}{Dζ and Fζ}}

\begin{definition}[$\Dz$ operator]
\label{def:dz}
\lean{FUST.DζOperator.Dζ}
\[
  \Dz[f](z) = \frac{\AFNum(\phiA\,f)(z) + \SymNum(\phiS\,f)(z)}{z},
\]
where $\AFNum$ and $\SymNum$ are 6-point convolutions over $\zs^k$ ($k=0,\ldots,5$).
The $\vp$-channels $\phiA$ and $\phiS$ are built from $\mathrm{Diff}_k$.
\end{definition}

\begin{definition}[$\Fz$ operator]
\label{def:fz}
\lean{FUST.FζOperator.Fζ}
\[
  \Fz[f](z) = \AFNum(5\,\phiA\,f)(z) + \SymNum(\phiS^{\mathrm{int}}\,f)(z),
\]
where $\phiS^{\mathrm{int}} = 5\,\phiS$ has $\GEI$-integral coefficients.
\end{definition}

\begin{theorem}
\label{thm:fz-eq}
\lean{FUST.FζOperator.Fζ_eq}
For $z\ne0$: $\Fz[f](z) = 5z\,\Dz[f](z)$.
\end{theorem}

\begin{theorem}[Gauge covariance]
\label{thm:gauge-covariance}
\lean{FUST.DζOperator.Dζ_gauge_covariance}
$\Dz[f(c\,\cdot\,)](z) = c\cdot\Dz[f](cz)$ for $c,z\ne 0$.
\end{theorem}

\section{\texorpdfstring{Kernel of $\Fz$}{Kernel of Fζ}}

\begin{theorem}[Mod~6 selection rule]
\label{thm:kernel-mod6}
\lean{FUST.FζOperator.Fζ_vanish_mod6_0,
  FUST.FζOperator.Fζ_vanish_mod6_2,
  FUST.FζOperator.Fζ_vanish_mod6_3,
  FUST.FζOperator.Fζ_vanish_mod6_4}
For all $k\in\N$:
\[
  \Fz[w^{6k}] = \Fz[w^{6k+2}] = \Fz[w^{6k+3}] = \Fz[w^{6k+4}] = 0.
\]
4 of every 6 monomial modes lie in $\ker(\Fz)$.
Active modes: $n\equiv 1,5\pmod6$.
\end{theorem}

\begin{theorem}[Eigenvalue formula]
\label{thm:eigenvalue}
\lean{FUST.FζOperator.Fζ_eigenvalue_mod6_1,
  FUST.FζOperator.Fζ_eigenvalue_mod6_5}
For $n\equiv1\pmod6$:
$\Fz[w^n](z) = (5\,c_A(n)\cdot\afcoeff + 6\,c_S(n))\,z^n$,
where $c_A(n)$, $c_S(n)$ are $\vp$-channel coefficients.
For $n\equiv5\pmod6$ the AF sign is negated.
\end{theorem}

\section{Hermitian Structure}

\begin{theorem}[AF anti-Hermitian, SY Hermitian]
\label{thm:hermitian}
\lean{FUST.FζOperator.AFNum_anti_hermitian,
  FUST.FζOperator.SymNum_hermitian}
$\AFNum(g)(\bar s) = -\overline{\AFNum(g)(s)}$ and
$\SymNum(g)(\bar s) = \overline{\SymNum(g)(s)}$
for Schwarz-reflective $g$.
On the real axis, $\AFNum$ is pure imaginary and $\SymNum$ is real.
\end{theorem}

\begin{theorem}[$\sqrt{15}$ factorization]
\label{thm:sqrt15}
\lean{FUST.FζOperator.eigenvalue_im_eq_sqrt15,
  FUST.FζOperator.eigenvalue_re_eq_phiS}
$\im(\lambda) = 10\sqrt{15}\,c_A$ (space),
$\re(\lambda) = 6\,c_S$ (time).
\end{theorem}

\section{Derivation Defect and Emergence}

\begin{definition}[Derivation defect]
\label{def:deriv-defect}
\lean{FUST.FζOperator.derivDefect}
\[
  \delta_{\Fz}(f,g)(z) = \Fz[fg](z) - f(z)\,\Fz[g](z) - \Fz[f](z)\,g(z).
\]
\end{definition}

\begin{theorem}[Emergence from kernel]
\label{thm:emergence}
\lean{FUST.FζOperator.emergence_matter,
  FUST.FζOperator.emergence_antimatter}
When $a,b$ are kernel modes but $a+b$ is active:
$\delta_{\Fz}(z^a, z^b) = \Fz[z^{a+b}]$.
Matter: $(6j+3)+(6k+4) \equiv 1\pmod6$.
Antimatter: $(6j+2)+(6k+3) \equiv 5\pmod6$.
\end{theorem}
